\sect{Probabilistic Relational Models}

% MLN in FOL and PL
So far we have studied \MarkovLogicNetworks{} in \propositionalLogic{} as probability distributions over worlds.
In FOL they define probability distributions over relations in worlds with a fixed set of objects.
More generally, such models are probabilistic relational models (see for an overview \cite{getoor_introduction_2019}.


We in this section treat random worlds in \firstOrderLogic{} with fixed domains $\worlddomain$.

%
We in this section show, when and how we can interpret likelihoods of \MarkovLogicNetworks{} in \firstOrderLogic{} in terms of samples of a \MarkovLogicNetwork{} in \propositionalLogic{}.

\subsect{\HybridFOLNetworks{}}

% Templates
Following \cite{richardson_markov_2006} \MarkovLogicNetworks{} in \firstOrderLogic{} are templates for distributions, which instantiate random worlds when choosing a set of objects $\worlddomain$.
Given a fixed set of constants, they then define a distribution over the worlds, which objects correspond with the constants. % this is database semantics!
This applies database semantics, where only those worlds are considered, where the unique name and domain closure assumptions given a set of constants are satisfied.

We count the number of true groundings to a formula by
\begin{align*}
    \countquantifier\enumfolformula|_{\worldindices}
    = \contraction{\groundingofat{\enumfolformula}{\indvariableof{\enumfolformula}}} \, .
\end{align*}
Using these counts as statistics, we now define \HybridFOLNetworks{} as a Computation Activation Network.

\begin{definition}[\HybridFOLNetworks{} (HFLN)]
    Let there be a set $\worlddomain$ of objects and a set $\folformulaset$ of \firstOrderLogic{} formulas.
    Further let $\hybridparam$ as in \defref{def:hybridLogicNetwork} be a tuple of a subset $\hardlegset\subset[\seldim]$, a tuple $\headindexof{\hardlegset}\in\bigtimes_{\selindex\in\hardlegset}[2]$ and $\canparamwithin$.
    We then define the \HybridFOLNetwork{} as the distribution
    \begin{align*}
        \probofat{\folhlnparameters}{\worldvariables}
        &= \breakablenormalizationof{\{\bencodingofat{\countquantifier\enumfolformula}{\headvariableof{\selindex},\worldvariables}\wcols\selindexin\} \\
        & \quad \cup \{\softactlegwith\wcols\selindexin\}
        \cup \{\kcoreofat{\selindex,\hardparam}{\headvariableof{\selindex}}\wcols\selindexin\}
        }{\worldvariables}
    \end{align*}
    where
    \begin{align*}
        \kcoreofat{\selindex,\hardparam}{\headvariableof{\selindex}}
        = \begin{cases}
              \fbasisat{\headvariableof{\selindex}} & \ifspace \selindex\in\hardlegset \ncond \headindexof{\selindex} = 0 \\
              \onehotmapofat{\headdimof{\selindex}-1}{\headvariableof{\selindex}} & \ifspace \selindex\in\hardlegset \ncond \headindexof{\selindex} = 1 \\
              \onesat{\headvariableof{\selindex}} & \text{else} \, .
        \end{cases}
    \end{align*}
\end{definition}



The mean parameter polytope is the convex hull of the vectors
\begin{align*}
    \sencodingofat{\folformulaset}{\indexedworldvariables,\selvariable}
\end{align*}
to the worlds $\worldindices$ with $\basemeasureat{\indexedworldvariables}=1$.
These vectors store are the counts of satisfied groundings to each formula, that is
\begin{align*}
    \sencodingofat{\folformulaset}{\indexedworldvariables,\selvariable} = \cardof{
        \{\indindexof{\enumfolformula} \wcols \groundingofat{\enumfolformula}{\indexedindvariableof{\enumfolformula}} = 1 \}
    } \, .
\end{align*}

\subsect{Propositionalization}

% Propositionalization
Let us notice, that different to the case of propositional \HybridLogicNetworks{} treated in \charef{cha:networkRepresentation}, the statistic does not consist of boolean features, when formulas contain variables and we have multiple objects.
One could, however, replace each $\enumfolformula$ by the set of the possible groundings, i.e. substitutions of the formulas variables by any tuple of objects in $\worlddomain$.
The resulting distribution would be an \HybridLogicNetwork{} with boolean statistic, which coincides with the \HybridFOLNetwork{} when posing certain weight sharing conditions on $\canparam$.
The downside of this construction is the increase in the number of features from $\seldim$ to $\sum_{\selindexin} \cardof{\worlddomain}^{\cardof{\indvariableof{\enumfolformula}}}$.
This polynomial in the cardinality of the domain set increase poses significant computational challenges, see \cite{richardson_markov_2006}.

We now show that the \HybridFOLNetwork{} can be propositionalized to a \HybridLogicNetwork{} in \propositionalLogic{}.

\begin{lemma}
    We have
    \begin{align*}
        &\contractionof{\bencodingofat{\countquantifier\enumfolformula}{\headvariableof{\selindex},\worldvariables},\softactlegwith}{\worldvariables} \\
        &\quad = \contractionof{
            \{\bencodingofat{\groundingofat{\enumfolformula}{\worldindices,\indindexof{\enumfolformula}}}{\headvariableof{\indindexof{\enumfolformula}}} \wcols \indindexofin{\enumfolformula}\}
            \cup
            \{\softactlegat{\headvariableof{\indindexof{\enumfolformula}}} \wcols \indindexofin{\enumfolformula}\}
        }{\worldvariables}
    \end{align*}
    and
    \begin{align*}
        &\contractionof{\bencodingofat{\countquantifier\enumfolformula}{\headvariableof{\selindex},\worldvariables},\kcoreofat{\selindex,\hardparam}{\headvariableof{\selindex}}}{\worldvariables} \\
        &\quad = \contractionof{
            \{\bencodingofat{\groundingofat{\enumfolformula}{\worldindices,\indindexof{\enumfolformula}}}{\headvariableof{\indindexof{\enumfolformula}}} \wcols \indindexofin{\enumfolformula}\}
            \cup
            \{\kcoreofat{\selindex,\hardparam}{\headvariableof{\indindexof{\enumfolformula}}} \wcols \indindexofin{\enumfolformula}\}
        }{\worldvariables} \, .
    \end{align*}
\end{lemma}
\begin{proof}
    The first claim holds, since for any world $\worldindices$ we have
    \begin{align*}
        &\contractionof{\bencodingofat{\countquantifier\enumfolformula}{\headvariableof{\selindex},\worldvariables},\softactlegwith}{\indexedworldvariables} \\
        &\quad = \prod_{\indindexof{\enumfolformula}\wcols\groundingofat{\enumfolformula}{\worldindices,\indindexof{\enumfolformula}}=1} \expof{\indexedcanparam} \\
        &\quad = \contractionof{
            \{\bencodingofat{\groundingofat{\enumfolformula}{\worldindices,\indindexof{\enumfolformula}}}{\headvariableof{\indindexof{\enumfolformula}}} \wcols \indindexofin{\enumfolformula}\}
            \cup
            \{\softactlegat{\headvariableof{\indindexof{\enumfolformula}}} \wcols \indindexofin{\enumfolformula}\}
        }{\indexedworldvariables} \, .
    \end{align*}
    For $\selindex\notin\hardlegset$ the second claim is trivial, since any contraction of a basis encoding with a trivial head is trivial.
    If $\selindex\in\hardlegset$ and $\headindexof{\selindex}=0$ then
    \begin{align*}
        &\contractionof{\bencodingofat{\countquantifier\enumfolformula}{\headvariableof{\selindex},\worldvariables},\kcoreofat{\selindex,\hardparam}{\headvariableof{\selindex}}}{\indexedworldvariables} \\
        &\quad=
        \begin{cases}
            1 & \ifspace \forall \indindexof{\enumfolformula} \defcols  \groundingofat{\enumfolformula}{\worldindices,\indindexof{\enumfolformula}} = 0\\
            0 & \text{else}
        \end{cases} \\
        &\quad= \contractionof{
            \{\bencodingofat{\groundingofat{\enumfolformula}{\worldindices,\indindexof{\enumfolformula}}}{\headvariableof{\indindexof{\enumfolformula}}} \wcols \indindexofin{\enumfolformula}\}
            \cup
            \{\kcoreofat{\selindex,\hardparam}{\headvariableof{\indindexof{\enumfolformula}}} \wcols \indindexofin{\enumfolformula}\}
        }{\indexedworldvariables} \, .
    \end{align*}
    For $\headindexof{\selindex}=1$ this holds analogously.
\end{proof}

Each substitution of the variables in $\enumfolformula$ by objects in $\worlddomain$, which satisfies the formula in the world $\worldindices$, therefore provides a factor of $\expof{\canparamat{\indexedselvariable}}$ to the probability of $\worldindices$.


%\red{Here we directly define them as exponential families distributing $\worldvariables$ for a given set of objects $\worlddomain$.}
%\red{To avoid a similar discussion as in \charef{cha:networkRepresentation} we directly allow for boolean base measures and call the distributions \HybridFOLNetworks{}.}
%
%\begin{definition}[\HybridFOLNetworks{} (HFLN)]
%    Let there be a set $\folformulaset$ of \firstOrderLogic{} formulas with maximal arity $\indorder$, which is enumerated by a selection variable $\selvariable$ of dimension $\seldim$.
%    Further, let there be a set of objects $\worlddomain$ and a boolean base measure $\basemeasureat{\shortindvariables}$.
%    The family of \HybridFOLNetworks{} $\expfamilyof{\restfolformulaset,\basemeasure}$ defined by the tuple $(\folformulaset,\worlddomain,\basemeasure)$ is the exponential family of joint distributions to the variables $\worldvariables$ with the statistics
%    \begin{align*}
%        \sstat^{\restfolformulaset}_{\selindex}\left[\indexedworldvariables\right]
%        = \contraction{\groundingof{\enumfolformula}}
%    \end{align*}
%    and the base measure $\basemeasure$.
%%    The \MarkovLogicNetwork{} instantiated for a given set of objects $\worlddomain$ and a base measure $\basemeasure$ is the random world, which is a member of the exponential family with sufficient statistics
%%    \begin{align*}
%%        \sstatcoordinateofat{\selindex}{\indexedworldvariables} = \contraction{\groundingof{\enumfolformula}}
%%        %\sstat_{\selindex}(\worldindices)  = \contraction{\groundingof{\folexformula_\selindex}} % Formulas can have different
%%    \end{align*}
%%    and canonical parameters $\weight$.
%\end{definition}
%
%Each element of the family $\expfamilyof{\restfolformulaset,\basemeasure}$ is represented by a canonical parameter $\canparamat{\selvariable}$.


%We will in the next sections explore an alternative way to apply the theory of \charef{cha:networkRepresentation} and \charef{cha:networkReasoning}, namely based on importance formulas.


%% Interpretation
%The statistics
%\begin{align*}
%    \contraction{\groundingof{\folexformula_\selindex}} % Formulas can have different
%\end{align*}
%can be interpreted as the number of substitutions to a formula, such that the formula ist satisfied.
%Each substitution satisfying a formula adds a factor of $\expof{\canparam_\selindex}$ to the probability of the respective world before normalization.


%
%When constructing a world tensor to a theory with predicates of different order, we already argued that we extend the arity of predicates by tensor products with $\onehotmapof{0}$.
%To define random world tensors, we then restrict the corresponding base measure to be supported only on those worlds where the extended predicates hold only at the individual $\exindividualof{0}$ at the extended axis.


% Comparison with PL MLN
%We choose extraction formulas $\extformulaof{\atomenumerator}$ such that any formula in the FOL MLN has a propositional equivalent (see \defref{def:propositionalEquivalent}).
%The statistic map is then a formula selecting tensor as in the propositional logic case contracted with the groundings of $\extformulaof{\atomenumerator}$.






\subsect{Conditioning on an Importance Formula}

%\red{Analogous to a guard formula in \cite[Definition 6.11]{koller_probabilistic_2009}!}

%The boolean base measure $\basemeasure$ of a \HybridFOLNetwork{} is the subset encoding of the possible worlds which have a non-vanishing probability with respect to any member of the family.
%We now construct specific base measures based on a fixed grounding tensor of an importance formula.
%This will reduce the number of object tuples influencing the probability distribution in order to arrive at an interpretation of FOL MLNs as likelihoods to datasets of propositional MLNs.

To reduce the number of object tuples influencing the probability distribution, we now condition \HybridFOLNetworks{} on situations where a formula, called the importance formula, has a fixed grounding tensor.

To this end, we mark pairs of term indices relevant to the distributions by an auxiliary index $\datindexin$.
Given a set $\{\indindexof{[\indorder]}^{\datindex} \wcols \datindexin \}$ of indices to the important tuples we build a set encoding (see \defref{def:subsetEncoding})
\begin{align*}
    \fixedimpformula = \sum_{\datindexin} \left(
    \bigotimes_{\indenumeratorin} \onehotmapofat{\indindexof{\indenumerator}^{\datindex}}{\indvariableof{\indenumerator}}
    \right) \, .
\end{align*}

% Interpretation as grounding
We interpret the tensor $\fixedimpformula$ as the grounding of a formula, which we call the importance formula.

% Restricting to worlds with identical grounding
To have a constant importance formula we define a syntactic representation and restrict the support of the \HybridFOLNetwork{} to those world coinciding with groundings of the importance formula coinciding with $\fixedimpformula$ by designing a base measure
\begin{align*}
    \fixedimpbm
    = \begin{cases}
          1 & \ifspace \groundingofat{\impformula}{\indvariableof{\impformula}} = \fixedimpformula \\
          0 & \text{else}
    \end{cases} \, .
\end{align*}

% Conditioning on exquery
The base measure restricts the \HybridFOLNetwork{} to be those worlds, where $\groundingof{\impformula}$ is coincides with the fixed tensor $\fixedimpformula$.
Intuitively, $\groundingof{\impformula}$ represents certain evidence about a \firstOrderLogic{} world, whereas other formulas are uncertain.


%\begin{assumption}
%    \label{ass:importanceBasemeasure}
%    Given a base measure $\fixedimpbm$, we assume that there is an importance formula $\impformulaat{\shortindvariables}$ such that
%    \begin{align*}
%        \fixedimpbm
%        = \begin{cases}
%              1 & \ifspace \groundingofat{\impformula}{\indvariableof{\impformula}} = \fixedimpformula \\
%              0 & \text{else}
%        \end{cases} \, .
%    \end{align*}
%\end{assumption}


\subsect{Decomposition of the Log Likelihood}


% Extraction query
To reduce the likelihood of a world to we make the assumption that all formulas in a \HybridFOLNetwork{} are of the form
\begin{align}
    \label{eq:folImplicationForm}
%    \folexformula_{\selindex}(\individuals) =
    \enumfolformulaat{\indvariableof{\enumfolformula}}
    = \left( \impformulaat{\shortindvariables} \Rightarrow \headfolformulaofat{\selindex}{\indvariableof{\enumfolformula}} \right)
\end{align}
that is a rule with the importance formula being the premise.
In particular, we assume, that they depend on all term variables $\shortindvariables$.
If this is not the case, we extend the formula trivially on the missing term variables.
When this assumption holds, we can think of the importance formula as a conditions on individuals to satisfy a statistical relation given by $\headfolexformula$.

Towards connecting with \propositionalLogic{}, we further make the assumption, that we can decompose the formula $\headfolformulaof{\selindex}$ in what we will call extraction formulas.

\begin{assumption}
    \label{ass:propositionalHeads}
    We assume that there exist formulas $\{\extformulaofat{\catenumerator}{\worldvariables,\shortindvariables} \wcols \catenumeratorin\}$, which we refer to as atom extraction formulas, and an importance formula $\impformulaat{\shortindvariables}$ such that the following holds.
    To each \firstOrderLogic{} formula $\enumfolformula$ there is another \firstOrderLogic{} formula $\headfolformulaofat{\selindex}{\indvariableof{\enumfolformula}}$ and a propositional formula $\enumformulaat{\shortcatvariables}$ such that
    \begin{align*}
        \enumfolformulaat{\worldvariables,\shortindvariables}
        = \left( \impformulaat{\worldvariables,\shortindvariables} \Rightarrow \headfolformulaofat{\selindex}{\worldvariables,\shortindvariables} \right)
    \end{align*}
    and
    \begin{align*}
        \headfolformulaofat{\selindex}{\shortindvariables} =
        \contractionof{
            \{\enumformulaat{\shortcatvariables}\} \cup \{\bencodingofat{\extformulaof{\catenumerator}}{\catvariableof{\catenumerator},\worldvariables,\shortindvariables}
            \wcols \catenumeratorin\}
        }{\shortindvariables} \, .
    \end{align*}
\end{assumption}

We depict the assumption, that any formula is of the form \eqref{eq:folImplicationForm} in the diagram
\begin{center}
    \input{tikz/distributions/implication_propositionalization.tex}
\end{center}
where the second summand depends only on the query $\impformula$ and therefore does not appear in the likelihood.


%\begin{example}[Trivial importance formula]
%	When the importance formula is always satisfied, any tuple of objects contributes to the likelihood.
%	This original approach to \MarkovLogicNetworks{} \cite{richardson_markov_2006} however leads to many datapoints which are also dependent on each other.
%\end{example}


\subsect{Reduction to Propositional Logic}

We now make additional assumptions to decompose the partition function of an \HybridFOLNetwork{} as a product of \HybridLogicNetwork{} partition functions.

\begin{assumption}
    \label{ass:independentTuples}
    Given an importance formula $\impformula$ and a tensor $\fixedimpformula$ we consider the base measure
    \begin{align*}
        \fixedimpbm
        = \begin{cases}
              1 & \ifspace \groundingofat{\impformula}{\indvariableof{\impformula}} = \fixedimpformulawith \\
              0 & \text{else}
        \end{cases}
    \end{align*}
    of worlds with grounding of $\impformula$ by $\fixedimpformula$.
    Let further enumerate by $\datindexin$ the support of $\fixedimpformula$, that is $\sampleind$ are the indices of $\shortindvariables$ with $\fixedimpformulaat{\shortindvariables=\sampleind}=1$.
    We assume that with respect to $\fixedimpformula$ the variables
    \begin{align}
        \left(\groundingofat{\extformulaof{\atomenumerator}}{\shortindvariables=\sampleind}\right)
    \end{align}
    are for $\atomenumeratorin$ and $\datindexin$ independent and uniformly distributed.
\end{assumption}

When the above assumption holds, we now show that the probability of a \firstOrderLogic{} world with respect to a \HybridFOLNetwork{} coincides with the likelihood of a dataset in a propositional \HybridLogicNetwork{}.

\begin{theorem}
    \label{the:FOLworldToPLdataset}
    Let there be a set of formulas $\folformulaset$ such that \assref{ass:propositionalHeads} and \assref{ass:independentTuples} hold with an importance formula $\impformula$ and a tensor $\fixedimpformula$.
    Then for any $\hybridparam$ and any world $\catindexof{\folworldsymbol}$ with $\groundingof{\impformula}=\fixedimpformula$ we have
    \begin{align*}
        \frac{1}{\datanum} \lnof{\condprobwrtof{\folhlnparameters}{\indexedworldvariables}{\groundingof{\impformula}=\fixedimpformula}}
        = \centropyof{\empdistribution}{\probof{\hlnparameters}} -
        \frac{\lnof{\contraction{\fixedimpbm}}}{\datanum}
    \end{align*}
    where $\formulaset$ is the set of propositional equivalents to $\folformulaset$ (see \assref{ass:propositionalHeads}) and $\datamap$ the data map with evaluation at $\datindexin$ by the enumerated non-vanishing coordinates of $\fixedimpformulawith$
    \begin{align*}
        \datapoint
        = \big(\groundingofat{\extformulaof{0}}{\shortindvariables=\sampleind},\ldots,\groundingofat{\extformulaof{\atomorder-1}}{\shortindvariables=\sampleind}\big) \, .
    \end{align*}
\end{theorem}

To show the theorem, we show first in the following lemma the factorization of the partition function of the \HybridFOLNetwork{}.

% NEW
\begin{lemma}
    \label{lem:FOLpartitionfunctionfactorization}
    Under the assumptions of \theref{the:FOLworldToPLdataset}, we have
    \begin{align*}
        &\contraction{
            \bencodingofat{\restfolformulaset}{\headvariables,\worldvariables},\hypercoreofat{\hybridparam}{\headvariables},\fixedimpbm
        } \\
        &\quad=\contraction{\fixedimpbm} \cdot
        \left(  \prod_{\shortindindices\wcols\fixedimpformulaat{\indexedshortindvariables}=0} \hypercoreofat{\selindex,\hybridparam}{\headvariable=1}\right) \cdot \\
        &\quad\quad \left(\frac{1}{2^{\atomorder}}
        \cdot \contraction{\hlnstatccwith,\hypercoreofat{\hybridparam}{\headvariables}}
        \right)^{\datanum} \, .
    \end{align*}
\end{lemma}
\begin{proof}
    Using the assumption on the structure of the formulas we have
    \begin{align*}
        \groundingof{\countquantifier\enumfolformula}
        = \groundingof{\countquantifier\lnot\impformula} + \groundingof{\countquantifier(\impformula\land\headfolformulaof{\selindex})}
    \end{align*}
    and for any $\worldindices$
    \begin{align*}
        &\contractionof{\bencodingofat{\countquantifier\enumfolformula}{\headvariableof{\selindex}},\hypercoreofat{\selindex,\hybridparam}{\headvariable},\fixedimpbm}{\indexedworldvariables} \\
        &\quad = \left(\prod_{\shortindindices\wcols\fixedimpformulaat{\indexedshortindvariables}=0} \hypercoreofat{\selindex,\hybridparam}{\headvariable=1}\right) \cdot
        \contractionof{\bencodingofat{\countquantifier(\impformula\land\headfolformulaof{\selindex})}{\headvariableof{\selindex}},\hypercoreofat{\selindex,\hybridparam}{\headvariableof{\selindex}},\fixedimpbm}{\indexedworldvariables} \\
    \end{align*}
    Here by $\hypercoreofat{\selindex,\hybridparam}{\headvariable}$ we denote the two-dimensional realization of the activation core.
    We use that the grounding tensor of $\impformula$ is constant among the by $\fixedimpbm$ supported worlds and thus
    \begin{align*}
        &\contractionof{\bencodingofat{\countquantifier(\impformula\land\headfolformulaof{\selindex})}{\headvariableof{\selindex}},\hypercoreofat{\selindex,\hybridparam}{\headvariableof{\selindex}},\fixedimpbm}{\worldvariables} \\
        % Using that impformula is constant
        &\quad =
        \contractionof{
            \bigcup_{\datindexin}\left\{\bencodingofat{\groundingofwrt{\headfolformulaof{\selindex}}{\sampleind}}{\headvariableof{\selindex,\datindex}},\hypercoreofat{\selindex,\hybridparam}{\headvariableof{\selindex,\datindex}}\right\}
            \cup\{\fixedimpbm\}
        }{\worldvariables} \, .
    \end{align*}
    For each $\datindexin$ we have by \assref{ass:propositionalHeads}
    \begin{align*}
        % Using that headformula decomposes to propositional
        &\contractionof{
            \bencodingofat{\groundingofwrt{\headfolformulaof{\selindex}}{\sampleind}}{\headvariableof{\selindex,\datindex}},\hypercoreofat{\selindex,\hybridparam}{\headvariableof{\selindex,\datindex}}
        }{\worldvariables} \\
        & \quad =
        \contractionof{
            \{\bencodingofat{\groundingofwrt{\extformulaof{\atomenumerator}}{\sampleind}}{\headvariableof{\atomenumerator,\datindex}} \wcols \atomenumeratorin \}
            \cup \{\bencodingofat{\enumformula}{\formulavar,\headvariableof{\atomenumerator,\datindex}},\hypercoreofat{\selindex,\hybridparam}{\formulavar}\}
        }{\worldvariables} \, .
    \end{align*}
    From \assref{ass:independentTuples} we know
    \begin{align*}
        \contractionof{
            \{\bencodingofat{\groundingofwrt{\extformulaof{\atomenumerator}}{\sampleind}}{\headvariableof{\atomenumerator,\datindex}}
            \wcols \atomenumeratorin \ncond \datindexin \} \cup \{\fixedimpbm\}
        }{\headvariableof{[\atomorder]\times[\datanum]}}
        = \frac{\contraction{\fixedimpbm}}{2^{\atomorder\cdot\datanum}} \cdot \onesat{\headvariableof{[\atomorder]\times[\datanum]}} \, .
    \end{align*}
    Combining the above we get
    \begin{align*}
        &\contraction{
            \bencodingofat{\restfolformulaset}{\headvariables,\worldvariables},\hypercoreofat{\hybridparam}{\headvariables},\fixedimpbm
        } \\
        &\quad=
        \frac{\contraction{\fixedimpbm}}{2^{\atomorder\cdot\datanum}} \cdot \\
        &\quad\quad \left(  \prod_{\shortindindices\wcols\fixedimpformulaat{\indexedshortindvariables}=0} \hypercoreofat{\selindex,\hybridparam}{\headvariable=1}\right)
        \prod_{\datindexin}\contraction{
            \bencodingofat{\hlnstat}{\headvariables,\headvariableof{[\atomorder]\times\{\datindex\}}},\hypercoreofat{\hybridparam}{\headvariables}
        } \\
        &\quad=\contraction{\fixedimpbm} \cdot
        \left(  \prod_{\shortindindices\wcols\fixedimpformulaat{\indexedshortindvariables}=0} \hypercoreofat{\selindex,\hybridparam}{\headvariable=1}\right) \\
        &\quad\quad \left(\frac{1}{2^{\atomorder}}
        \cdot \contraction{\hlnstatccwith,\hypercoreofat{\hybridparam}{\headvariables}}
        \right)^{\datanum} \, . \qedhere
    \end{align*}
\end{proof}

% For
We notice that for the event $\groundingof{\impformula}=\fixedimpformula$ to be of non-vanishing probability, we need to have $\headindexof{\hardlegset}=\ones_\hardlegset$ or $\fixedimpformulawith=\onesat{\indvariableof{\impformula}}$.
With this lemma, we are now show \theref{the:FOLworldToPLdataset}.

\begin{proof}[Proof of \theref{the:FOLworldToPLdataset}]
    We have
    \begin{align*}
        \lnof{\condprobwrtof{\folhlnparameters}{\indexedworldvariables}{\groundingof{\impformula}=\fixedimpformula}}
        &= \lnof{\contractionof{
            \bencodingofat{\restfolformulaset}{\headvariables,\worldvariables},\hypercoreofat{\hybridparam}{\headvariables},\fixedimpbm
        }{\indexedworldvariables}} \\
        &\quad - \lnof{\contraction{
            \bencodingofat{\restfolformulaset}{\headvariables,\worldvariables},\hypercoreofat{\hybridparam}{\headvariables},\fixedimpbm
        }}
    \end{align*}
    While the second term is decomposed by \lemref{lem:FOLpartitionfunctionfactorization} we now derive a decomposition of the first term.
    By \assref{ass:propositionalHeads} and $\groundingof{\impformula}=\fixedimpformula$ we have
    \begin{align*}
        &\contractionof{
            \bencodingofat{\restfolformulaset}{\headvariables,\worldvariables},\hypercoreofat{\hybridparam}{\headvariables},\fixedimpbm
        }{\indexedworldvariables} \\
        &\quad =
        \left(\prod_{\shortindindices\wcols\fixedimpformulaat{\indexedshortindvariables}=0} \hypercoreofat{\selindex,\hybridparam}{\headvariable=1}\right)
        \cdot \\
        &\quad\quad \prod_{\datindexin}
        \contraction{
            \{\bencodingofat{\groundingofwrt{\extformulaof{\atomenumerator}}{\sampleind}}{\headvariableof{\atomenumerator},\indexedworldvariables} \wcols \atomenumeratorin \}
            \cup \{\hlnstatccwith,\hypercoreof{\hybridparam}{\headvariables}\}
        } \, .
    \end{align*}
    With \lemref{lem:FOLpartitionfunctionfactorization} we then have
    \begin{align*}
        &\frac{1}{\datanum}\lnof{\condprobwrtof{\folhlnparameters}{\indexedworldvariables}{\groundingof{\impformula}=\fixedimpformula}} \\
        &\quad= \frac{1}{\datanum}\sum_{\datindexin} \lnof{\contraction{
            \{\bencodingofat{\groundingofwrt{\extformulaof{\atomenumerator}}{\sampleind}}{\headvariableof{\atomenumerator},\indexedworldvariables} \wcols \atomenumeratorin \}
            \cup \{\hlnstatccwith,\hypercoreof{\hybridparam}{\headvariables}\}
        }} \\
        &\quad\quad - \lnof{\contraction{
            \{\bencodingofat{\groundingofwrt{\extformulaof{\atomenumerator}}{\sampleind}}{\headvariableof{\atomenumerator},\worldvariables} \wcols \atomenumeratorin \}
            \cup \{\hlnstatccwith,\hypercoreof{\hybridparam}{\headvariables}\}
        }} - \frac{\lnof{\contraction{\fixedimpbm}}}{\datanum} \\
        &\quad = \centropyof{\empdistribution}{\probof{\hlnparameters}} - \frac{\lnof{\contraction{\fixedimpbm}}}{\datanum} \, . \qedhere
    \end{align*}
\end{proof}

% Independent data investigation
Let us now investigate, in which cases the \assref{ass:independentTuples} of independent data can be matched.

\begin{example}
    If the $\impformula$ and $\extformulas$ are predicates applied on variables, and the index tuples $\sampleind$ are pairwise different, then \assref{ass:independentTuples} is met.
    This is the case, since the values $\groundingof{\extformulaofat{\atomenumerator}{\shortindvariables=\sampleind}}$ are determined by different variables in $\worldvariables$.
\end{example}



There are situations, where \assref{ass:independentTuples} is violated.
\begin{itemize}
    \item extraction formula being a) conjunctions of predicates: Probability that they are satisfied decreases
    b) disjunctions of predicates: Probability that they are satisfied increases
    \item extraction formula coinciding with importance formula: Always satisfied, in this case still boolean
    \item extraction formulas contradicting each other, more general not independent from each other
\end{itemize}

%Let us notice, that non-boolean base measures could be treated in a same manner, but several developments in this work, such as cross-entropy decompositions in \charef{cha:probReasoning} would receive further terms.



\begin{remark}[Approximation by Independent Samples]
    As argued above, we do not have independent samples in general.
    As a consequence, we cannot apply \lemref{lem:FOLpartitionfunctionfactorization} to decompose the partition function term of the log-probability into factors to each solution map of $\impformula$.
    In this case, it might be still benefitial to use the reduction to the likelihood of a HLN, but needs to understand it as a approximation to the true world probability.

    %
    If the expectations of each sample with respect to the marginalized distributions coincide, the average of empirical distribution also coincides with these (by linearity).
    When the creation of samples has sufficient mixing properties, the empirical distribution converges to this expectation in the asymptotic case of large numbers of samples.

\end{remark}



\subsect{Sample Extraction from First-Order Logic Worlds}

We have observed that in certain situations the log-likelihood of a \firstOrderLogic{} world with respect to a \HybridFOLNetwork{} coincides with the likelihood of a data set in a propositional \HybridLogicNetwork{}.
Let us now investigate the extraction process of these data set.
%The decomposition of the likelihood suggests the following approach to generate samples from groundings:
%%We propose the following approach to generate datacores from groundings:
%\begin{itemize}
%    \item Define an importance formula $\impformula$, which we decompose in the basis CP decomposition and interpret each slice as the one-hot encoding of the datapoint.
%    \item Define for $\atomenumeratorin$ extraction formulas $\extformulaof{\atomenumerator}$ generating the atoms $\catvariableof{\atomenumerator}$.
%    %Predicates along with assignment of variables / constants to its positions.
%    \item Contract the groundings of each formula $\extformulaof{\atomenumerator}$ with the grounding of $\impformula$ to build a data core.
%\end{itemize}
%\subsubsect{Representation by Tensor Networks}
We model the extraction process as a relation between a tuple of individuals and the extracted world in the factored system of atoms $\catvariableof{\atomenumerator}$.

\begin{definition}
    \label{def:extractionRelation}
    Given a \firstOrderLogic{} world $\worldindices$, an importance formula $\impformula$ and extraction formulas $\extformulaof{\catenumerator}$ for $\catenumeratorin$, we define the extraction relation
    \begin{align*}
        \extractionrelation \subset \left(\symindstates\right) \otimes \left(\atomstates\right)
    \end{align*}
    by
    \begin{align*}
        \extractionrelation
        = \{ (\shortindindices, \shortcatindices)
        \wcols  \groundingofat{\impformula}{\indexedshortindvariables} = 1 \ncond \uniquantwrtof{\catenumeratorin}{\catindexof{\atomenumerator} = \extformulaofat{\atomenumerator}{\indexedshortindvariables}} \} \, .
    \end{align*}
\end{definition}

The encoding of an extraction relation is the tensor
\begin{align*}
    \bencodingofat{\extractionrelation}{\shortindvariables,\shortcatvariables} \subset \left(\indspace\right) \otimes \left(\atomspace\right) \,
\end{align*}
and drawn in a contraction diagram by
\begin{center}
    \input{tikz/distributions/extraction_relation.tex}
\end{center}
Here the contraction of $\bencodingof{\impformula}$ with the truth vector $\tbasis$ represents the matching condition posed by $\impformula$ when extracting pairs of individuals.

%% Empirical Distribution
The empirical distribution of the extracted data is then the normalized contraction leaving only the legs to the extracted atomic formulas open, that is
\begin{align*}
    \empdistribution
    = \frac{
        \contractionof{\bencodingof{\extractionrelation}}{\shortcatvariables}
    }{
        \contraction{\bencodingof{\extractionrelation}}
    }  \, .
\end{align*}
Here the number of extracted data is the denominator
\begin{align*}
    \datanum
    = \contraction{\bencodingof{\extractionrelation}}
    = \contraction{\bencodingofat{\impformula}{\headvariableof{\impformula},\shortindvariables},\tbasisat{\headvariableof{\impformula}}}\, .
\end{align*}

We depict this by
\begin{center}
    \input{tikz/distributions/empirical_generation.tex}
\end{center}

%\subsubsect{Decomposition of extracted data}

To connect with the empirical distribution introduced in \secref{sec:empDistribution} we now show how the empirical distribution extracted from the interpretations of the formulas $\impformula,\extformulas$ on a \firstOrderLogic{} world $\worldindices$ can be represented by tensor networks.

First of all, we decompose the importance formula into a basis $\cpformat$ format (see \charef{cha:sparseRepresentation}), that is a decomposition
\begin{align*}
    \groundingofat{\impformula}{\shortindvariables}
    = \contractionof{
        \{\legcoreofat{\indenumerator}{\indvariableof{\indenumerator},\datvariable} \wcols \indenumeratorin \}
    }{\shortindvariables}
\end{align*}
such that all $\legcoreofat{\indenumerator}{\indvariableof{\indenumerator},\datvariable}$ are directed and boolean tensors.
Here an auxiliary variables $\datvariable$ taking values in $[\datanum]$ is introduced, which we call the data variable, which enumerates the non-vanishing coordinates of $\groundingof{\impformula}$.
With this decomposition, we can understand the decomposition of $\groundingofat{\impformula}{\shortindvariables}$ as a basis encoding of an term selection map $\secdatamap$ with coordinate maps defined such that
\begin{align*}
    \bencodingofat{\secdatamap_{\indenumerator}}{\indvariableof{\indenumerator},\datvariable}
    = \legcoreofat{\indenumerator}{\indvariableof{\indenumerator},\datvariable} \, .
\end{align*}
We depict this decomposition by:
\begin{center}
    \input{tikz/distributions/impformula_cp.tex}
\end{center}

Based on these construction, we now provide a tensor network decomposition of the extracted empirical distribution.

\begin{theorem}
    \label{the:extractionrelationDecomposition}
    Given a \firstOrderLogic{} world $\worldindices$, an importance formula $\impformula$ and extraction formulas $\extformulaof{\catenumerator}$ for $\catenumeratorin$, we have
    \begin{align*}
        \bencodingofat{\extractionrelation}{\shortindvariables,\shortcatvariables} =
        \contractionof{
            \{\bencodingofat{\groundingof{\extformulaof{\atomenumerator}}}{\catvariableof{\catenumerator},\shortindvariables} \wcols \catenumeratorin\}
            \cup \{\bencodingofat{\secdatamap_{\indenumerator}}{\indvariableof{\indenumerator},\datvariable} \wcols \indenumeratorin\}
        }{\shortindvariables,\shortcatvariables}
    \end{align*}
    and thus
    \begin{align*}
        \empdistributionat{\shortcatvariables} =
        \frac{1}{\datanum}  \contractionof{
            \{\bencodingofat{\groundingof{\extformulaof{\atomenumerator}}}{\catvariableof{\catenumerator},\shortindvariables} \wcols \catenumeratorin\}
            \cup \{\bencodingofat{\secdatamap_{\indenumerator}}{\indvariableof{\indenumerator},\datvariable} \wcols \indenumeratorin\}
        }{\shortcatvariables} \, .
    \end{align*}
\end{theorem}
\begin{proof}
    To show the first claim, let us choose arbitrary state tuples $\shortindindices$ and $\shortcatindices$.
    We then have
    \begin{align*}
        &\contractionof{
            \{\bencodingofat{\groundingof{\extformulaof{\atomenumerator}}}{\catvariableof{\catenumerator},\shortindvariables} \wcols \catenumeratorin\}
            \cup \{\bencodingofat{\secdatamap_{\indenumerator}}{\indvariableof{\indenumerator},\datvariable} \wcols \indenumeratorin\}
        }{\indexedshortindvariables,\indexedshortcatvariables} \\
        & \quad  =  \contraction{
            \{\bencodingofat{\groundingof{\extformulaof{\atomenumerator}}}{\indexedcatvariableof{\catenumerator},\indexedshortindvariables} \wcols \catenumeratorin\}
            \cup \{\bencodingofat{\secdatamap_{\indenumerator}}{\indexedindvariableof{\indenumerator},\datvariable} \wcols \indenumeratorin\}
        } \, .
    \end{align*}
    This contraction evaluates to $1$, if and only if for all $\catenumeratorin$ we have $\bencodingofat{\groundingof{\extformulaof{\atomenumerator}}}{\catvariableof{\catenumerator},\shortindvariables}=1$ and
    \begin{align*}
        \contraction{\{\bencodingofat{\secdatamap_{\indenumerator}}{\indexedindvariableof{\indenumerator},\datvariable} \wcols \indenumeratorin\}}  = 1 \, .
    \end{align*}
    The first condition is equal to $\catindexof{\atomenumerator} = \extformulaofat{\atomenumerator}{\indexedshortindvariables}$ for all $\catenumeratorin$ and the second to
    \begin{align*}
        \groundingofat{\impformula}{\indexedshortindvariables} = 1 \, .
    \end{align*}
    Comparing with the definition of the extraction relation (see \defref{def:extractionRelation}), we notice that these conditions are equal to $(\shortindindices,\shortcatindices)\in\extractionrelation$ and therefore to
    \begin{align*}
        \bencodingofat{\extractionrelation}{\indexedshortindvariables,\indexedshortcatvariables} \, .
    \end{align*}
    The first claim follows, since $\bencodingof{\extractionrelation}$ is boolean, as is the contraction of the cores $\bencodingof{\groundingof{\extformulaof{\atomenumerator}}}$ with the cores $\bencodingof{\secdatamap_{\indenumerator}}$, which leaves the outgoing variables $\shortcatvariables$ open.
    The second claim follows from the first using that $\empdistributionat{\shortcatvariables}=\frac{1}{\datanum}\contractionof{\bencodingof{\extractionrelation}}{\shortcatvariables}$.
\end{proof}

To connect with the representation of empirical distributions based on data cores (see \secref{sec:empDistribution}), we now form data cores by contractions with the grounding of extraction formulas with the cores $\bencodingof{\secdatamap_{\indenumerator}}$ (see \figref{fig:datacoreGeneration}),
\begin{align*}
    \datacoreofat{\atomenumerator}{\catvariableof{\catenumerator},\datvariable}
    = \contractionof{
        \{\bencodingofat{\groundingof{\extformulaof{\atomenumerator}}}{\catvariableof{\catenumerator},\shortindvariables}\}
        \cup \{ \legcoreofat{\indenumerator}{\indvariableof{\indenumerator},\datvariable} \wcols \indenumeratorin\}
    }{\catvariableof{\atomenumerator},\datvariable} \, .
\end{align*}

% Empirical distribution
The empirical distribution is then a tensor network of these tensors, as we show next.

\begin{theorem}
    \label{the:extractionDataCores}
    We have
    \begin{align*}
        \contractionof{\bencodingof{\extractionrelation}}{\shortcatvariables}
        = \contractionof{\{\datacoreofat{\atomenumerator}{\datvariable,\catvariableof{\atomenumerator}} \wcols \atomenumeratorin\}}{\shortcatvariables}
    \end{align*}
    and thus
    \begin{align*}
        \empdistributionat{\shortcatvariables}
        = \frac{1}{\datanum} \contractionof{\{\datacoreofat{\atomenumerator}{\datvariable,\catvariableof{\atomenumerator}}  \wcols \atomenumeratorin\}}{\shortcatvariables} \, .
    \end{align*}
\end{theorem}
\begin{proof}
    By \theref{the:extractionrelationDecomposition} we have
    \begin{align*}
        \bencodingofat{\extractionrelation}{\shortindvariables,\shortcatvariables} =
        \contractionof{
            \{\bencodingofat{\groundingof{\extformulaof{\atomenumerator}}}{\catvariableof{\catenumerator},\shortindvariables} \wcols \catenumeratorin\}
            \cup \{\bencodingofat{\secdatamap_{\indenumerator}}{\indvariableof{\indenumerator},\datvariable} \wcols \indenumeratorin\}
        }{\shortindvariables,\shortcatvariables} \, .
    \end{align*}
    Since $\bencodingofat{\secdatamap_{\indenumerator}}{\indvariableof{\indenumerator},\datvariable}$ are directed and boolean, they can be copied and separately contracted with each $\groundingof{\extformulaof{\atomenumerator}}$, without changing the contraction.
    We arrive at
    \begin{align*}
        &\bencodingofat{\extractionrelation}{\shortindvariables,\shortcatvariables} \\
        &\quad = \contractionof{
            \big\{\contractionof{
                \{\bencodingofat{\groundingof{\extformulaof{\atomenumerator}}}{\catvariableof{\catenumerator},\shortindvariables}\}
                \cup \{\bencodingofat{\secdatamap_{\indenumerator}}{\indvariableof{\indenumerator},\datvariable} \wcols \indenumeratorin\}
            }{\catvariableof{\catenumerator},\datvariable} \wcols \catenumeratorin \big\}
        }{\shortindvariables,\shortcatvariables} \\
        & \quad =  \contractionof{\{\datacoreofat{\atomenumerator}{\datvariable,\catvariableof{\atomenumerator}}  \wcols \atomenumeratorin\}}{\shortcatvariables} \, ,
    \end{align*}
    which established the claim.
\end{proof}

% Efficient contraction: Do also basis decomposition of the extraction query and use efficient contraction!
%Towards efficient calculation of the data cores, we build a basis CP decomposition of $\groundingof{\impformula}$, where we further demand $\scalarcore=\ones$.
%This is a collection of basis leg cores $\legcoreof{\fixedimpformula,\indenumerator}$ such that
%\begin{align*}
%    \fixedimpformula[\shortindvariablelist]
%    = \contractionof{ \left\{ \legcoreofat{\fixedimpformula,\indenumerator}{\datvariable,\indvariableof{\indenumerator}} \wcols \indenumeratorin \right\} }{\shortindvariablelist} \, .
%\end{align*}

% Data enumeration -> To representation
%We can further utilize any decomposition of $\impformula$ into a directed and binary CP Format to enumerate the datapoints by the slice index $\datindex$. % Approaches like SPARQL directly give us these by solution mappings.
%Understanding $\impformula$ as a query on the world being the database, such decomposition is given by the set of solution mappings.


\begin{figure}[t]
    \begin{center}
        \input{tikz/distributions/datacore_generation.tex}
    \end{center}
    \caption{Generation of a data core for the variable $\catvariableof{\catenumerator}$ given an extraction formula $\extformulaof{\catenumerator}$ and an importance formula, which grounding is decomposed into a basis CP format with leg vectors $\bencodingofat{\secdatamap_{\indenumerator}}{\indvariableof{\indenumerator},\datvariable}$.
    Term variables, which are appearing in the importance formula, but not in the extraction formula $\extformulaof{\catenumerator}$ can be treated trivally by contraction with the trivial tensor (here $\indvariableof{4},\ldots,\indvariableof{\indorder-1})$.
    }
    \label{fig:datacoreGeneration}
\end{figure}


% Comment: Exploitation of common structure
When many atom extraction formulas differ only by a constant, we can replace the constant by an auxiliary term variable.
The atoms are then the atomizations of this variable (see \secref{sec:categoricalTN}), treated as a categorical variable, with respect to the constant in the extraction query.
The advantages are that we can avoid the $\bencodingof{}$-formalism and directly model the categorical distributions.

This also enables a batchwise computation of multiple $\sparql$ queries, which differ only in one constant.


%\subsect{Design of the Formulas}
%
%Most intuitive when labeling individuals by classes.
%Extraction formulas $\extformulas$ can then be defined by subclasses of the member of a class and relations between objects of different classes. % Koller calls atomic formulas the template attributes
%We then choose $\formulaset$ as more involved formulas decomposed into connectives acting on these atoms.
%The importance formula $\impformula$ is then designed based on class memberships to ensure, that the arguments of the formulas are always of specific classes. % Koller specifies to each argument of the attributes a class
%
%% Approach
%We propose to
%\begin{itemize}
%    \item Execute an extraction query to get pairs of individuals (the pairDf).
%    \item Propositionalize the FOL Formulas independently on each tuple taking the individuals as a set of constant and filtering on the possible properties of each individuals.
%    (Can understand as adding knowledge that most of the relations do not hold)
%    \item Understand each such generated knowledge base as datapoint and average over them to get the empirical distribution to be fit.
%    \item Fit a MLN describing the statistical relations of unseen results of the extraction query, based on likelihood maximation.
%\end{itemize}




\sect{Generation of First-Order Logic Worlds}

\red{
    So far we have discussed, how Probabilistic Relational Models for \firstOrderLogic{} Knowledge Bases such as Knowledge Graphs can be built by extracting data.
    Conversely, any binary tensor can be interpreted as a Knowledge Graph.
    To be more precise, we follow the intuition that the ones coordinates mark possible worlds compatible with the knowledge about a factored system.
    Each possible world can then be encoded in a subgraph of the Knowledge Graph representing the world.
%
    This amounts to an "inversion" of the data generation process described in the subsection above.
}

In the previous section we have described a way to extract an effective empirical distribution for the likelihood of a \firstOrderLogic{} world given a \HybridFOLNetwork{}.
We now want to investigate methods to reproduce an empirical distribution based on a constructed \firstOrderLogic{} world.

\begin{definition}[Reproduction of Empirical Distributions]
    Given an empirical distribution $\empdistribution\in\atomspace$, we say that a triple $(\worldindices,\impformula,\shortextformulas)$ of a \firstOrderLogic{} world $\worldindices$ an importance formula $\impformula$ and extraction formulas $\shortextformulas=\{\extformulaof{\atomenumerator}\,:\,\atomenumeratorin\}$ reproduces $\empdistribution$, when
    \begin{align*}
        \empdistribution
        = \normalizationof{\{\groundingofat{\impformula}{\shortindvariables}\}\cup
        \{\bencodingofat{\kggroundingof{\extformulaof{\atomenumerator}}}{\catvariableof{\catenumerator},\shortindvariables} \wcols \atomenumeratorin\}
        }{\shortcatvariables} \, .
    \end{align*}
\end{definition}

% If \datamap is not known
Note that for distribution $\probtensor$ to be reproducible, it needs to have rational coordinates. %, since each coordinate can be interpreted as the frequency of the respective world in the data $\datamap$.
If any only if all coordinates are rational, we find a $\datanum\in\nn$ such that $\imageof{\datanum\cdot\probtensor}\subset\nn$.
We can then interpret $\datanum$ as the number of samples, and construct a sample selector map by understanding each coordinate of $\datanum\cdot\probtensor$ as the number of appearances of the respective world in the samples.

We show different schemes and give examples on Knowledge Graphs, where we provide examples for importance and extraction formulas by $\sparql$ queries.


%\subsect{Example: Generation of Knowledge Graphs} % To generation of \firstOrderLogic{} worlds?
%
% Having a directed and binary CP decomposition of $\exformula$, each possible world is encoded by a slice.


% Formalization
%\begin{definition}[Reproduction of Empirical Distributions]
%    Given an empirical distribution $\empdistribution\in\bigotimes_{\atomenumeratorin}\rr^2$, we say that a tuple $(\kg,\impformula,\{\extformulas\})$ of a Knowledge Graph $\kg$ and queries $\impformula,\extformulaof{\atomenumerator}$ reproduces $\empdistribution$, when
%    \[\empdistribution = \normalizationof{\{\kggroundingof{\impformula}\}\cup\{\bencodingof{\kggroundingof{\extformulaof{\atomenumerator}}\wcols \atomenumeratorin}\}}{\shortcatvariables} \, .  \]
%\end{definition}

%

%In a frequentist interpretation we instantiate each world according to the rate $\probtensor(\atomindices)$.
%This interpretation requires a rounding of the real probabilities by rational numbers.


\subsect{Samples by Single Objects}

%\subsect{Samples by single objects}

In the first reproduction scheme we construct datapoints by dedicated objects, which represent a sample, that is we choose a domain $\worlddomain=[\datdim]$.

\begin{theorem}
    \label{the:reproducingSingleObjects}
    Let there be an empirical distribution $\empdistribution$ to a sample selector map $\datamap$ (see \defref{def:dataMap}), we construct a world $\worldindices[\selvariable,\indvariable]$ with $\atomorder$ unary predicates by
    \begin{align*}
        \worldindices[{\selvariable,\indvariable}]
        = \sum_{\atomenumeratorin} \sum_{\datindexin \wcols \datamap_{\atomenumerator}(\datindex)=1} \onehotmapofat{\atomenumerator}{\selvariable} \otimes \onehotmapofat{\datindex}{\indvariable} \, .
    \end{align*}
    We further choose a trivial importance query, that is
    \begin{align*}
        \groundingofat{\impformula}{\indvariable} = \onesat{\indvariable} \, ,
    \end{align*}
    and extraction queries coinciding with the unary predicates, that is for $\atomenumeratorin$
    \begin{align*}
        \extformulaof{\atomenumerator} = \folpredicateof{\atomenumerator} \, .
    \end{align*}
    Then, the triple $(\worldindices,\impformula,\shortextformulas)$ reproduces $\empdistribution$.
%    reproduces with the trivial importance query and extraction queries coinciding with the predicates the dataset $\datamap$.
\end{theorem}
\begin{proof}
    By \theref{the:extractionDataCores} it is enough to show, that the data cores constructed from the data extraction process coincide with those of $\empdistribution$.
    We enumerate to this end the non-vanishing coordinates of $\groundingof{\impformula}$ by the data variable $\datvariable$ taking values $\datindexin$, as
    \begin{align*}
        \groundingofat{\impformula}{\indvariable=\datindex} = 1 \,
    \end{align*}
    and choose
    \begin{align*}
        \secdatamap = \identity \, .
    \end{align*}
    For arbitrary $\atomenumeratorin$ and $\datindexin$ we now have
    \begin{align*}
        \datacoreofat{\atomenumerator}{\catvariableof{\catenumerator},\indexeddatvariable}
        &= \contractionof{
            \bencodingofat{\groundingof{\extformulaof{\atomenumerator}}}{\catvariableof{\catenumerator},\indvariable},
            \legcoreofat{0}{\indvariable,\indexeddatvariable}
        }{\catvariableof{\atomenumerator},\datvariable} \\
        &= \contractionof{
            \bencodingofat{\groundingof{\extformulaof{\atomenumerator}}}{\catvariableof{\catenumerator},\indvariable},
            \onehotmapofat{\secdatamap(\datindex)}{\indvariable}
        }{\catvariableof{\atomenumerator},\indexeddatvariable} \\
        &= \onehotmapofat{\datamap_\atomenumerator(\datindex)}{\catvariableof{\catenumerator}} \, .
    \end{align*}
    This coincides with the slice of the data core of the CP representation of empirical distributions used in \theref{the:empCPRep}.
    Since the slice and the core was arbitrary, the tensor network representations in \theref{the:empCPRep} and \theref{the:extractionDataCores} are equal and thus the triple $(\worldindices,\impformula,\shortextformulas)$ reproduces $\empdistribution$.
\end{proof}


We now give by the next theorem an example of a Knowledge Graph with $\sparql$ queries reproducing and arbitrary empirical distribution.

\begin{theorem}
    \label{the:reproducingKGSingelObjects}
    Let $\empdistribution$ be an empirical distribution to the sample selector $\datamap$.
    We construct a Knowledge Graph of the resources $\worlddomain = \{s_\datindex \wcols \datindexin\} \cup \{C\} \cup \{C_\atomenumerator \wcols \atomenumeratorin\}$, where $s_{\datindex}$ represent samples and $C_\atomenumerator$ unary predicates, by
    \begin{align*}
        \kggroundingof{\rdf}
        =
        \sum_{\datindexin}
        \onehotmapof{\indexinterpretationof{s_\datindex}}{\sindvariable}
        \otimes \onehotmapof{\indexinterpretationof{\mathrdftype}}{\pindvariable}
        \otimes \onehotmapof{\indexinterpretationof{C}}{\oindvariable}
        +
        \sum_{\datindexin} \sum_{\atomenumeratorin \wcols \datamap_{\atomenumerator}(\datindex)=1}
        \onehotmapof{\indexinterpretationof{s_\datindex}}{\sindvariable}
        \otimes \onehotmapof{\indexinterpretationof{\mathrdftype}}{\pindvariable}
        \otimes \onehotmapof{\indexinterpretationof{C_\atomenumerator}}{\oindvariable} \, .
    \end{align*}
    We further define an importance formula by the $\sparql$ query
    \begin{centeredscript}
        \impformula = SELECT \{ ?x \} WHERE \{ ?x \quad \rdftype\quad C \, .\}
    \end{centeredscript}
    and for each $\atomenumeratorin$ an extraction formula by the query
    \begin{centeredscript}
        $\extformulaof{\atomenumerator}$ = SELECT \{ ?x \} WHERE \{ ?x \quad \rdftype \quad $C_\atomenumerator$ \, .\} \, .
    \end{centeredscript}
    Then the triple $(\kg,\impformula,\shortextformulas)$ reproduces $\empdistribution$.
\end{theorem}
\begin{proof}
    We show the theorem analogously to \theref{the:reproducingSingleObjects}, with the slide difference in the importance formula.
    We have for the grounding of $\impformula$ on $\kg$ that
    \begin{align*}
        \kggroundingofat{\impformula}{\indvariable} = \sum_{\datindexin}  \onehotmapof{\indexinterpretationof{s_\datindex}}{\indvariable}
    \end{align*}
    and enumerate the non-vanishing coordinates by $\datvariable$.

    For each extraction formula we have
    \begin{align*}
        \kggroundingofat{\extformulaof{\atomenumerator}}{\indvariable} = \sum_{\datindexin \wcols \datamap_{\atomenumerator}(\datindex)=1} \onehotmapof{\indexinterpretationof{s_\datindex}}{\indvariable} \,.
    \end{align*}
    It follows that the data cores used in \theref{the:extractionDataCores} are
    \begin{align*}
        \bencodingofat{\datamap_\atomenumerator}{\catvariableof{\atomenumerator},\datindex}
        = \onehotmapofat{0}{\catvariableof{\atomenumerator}} \otimes \left(\sum_{\datindexin \, : \, \datamap_{\atomenumerator}(\datindex)=0} \onehotmapofat{\datindex}{\datvariable}\right)
        +\tbasisat{\catvariableof{\atomenumerator}} \otimes \left(\sum_{\datindexin \, : \, \datamap_{\atomenumerator}(\datindex)=1} \onehotmapofat{\datindex}{\datvariable}\right)
    \end{align*}
    and they thus coincide with those in the decomposition in \theref{the:empCPRep}.
    The claim follows therefore with the same argumentation as in the proof of \theref{the:reproducingSingleObjects}.
\end{proof}

%
Let us provide some more insights on the construction of the reproducing Knowledge Graph in \theref{the:reproducingKGSingelObjects}.
By the insertions to the one-hot encodings $\onehotmapof{\indexinterpretationof{s_\datindex}}{\sindvariable} \otimes \onehotmapof{\indexinterpretationof{\mathrdftype}}{\pindvariable} \otimes \onehotmapof{\indexinterpretationof{C}}{\oindvariable}$ we mark each sample representing resource by a class and ensure its appearance as a $\mathrm{owl:NamedIndividual}$ in the graph.
The insertions $\onehotmapof{\indexinterpretationof{s_\datindex}}{\sindvariable}\otimes \onehotmapof{\indexinterpretationof{\mathrdftype}}{\pindvariable} \otimes \onehotmapof{\indexinterpretationof{C_\atomenumerator}}{\oindvariable}$ on the other side encode the sample selecting map, by inserting exactly the assertions corresponding with the respective sample.
%
In this simple Knowledge Graph, Description Logic is expressive enough to represent any formula $\folexformula$ composed of the formulas $\extformulas$.

%
%\begin{theorem}
%    Let there any empirical distribution $\empdistribution\in\bigotimes_{\atomenumeratorin}\rr^2$ and $\datanum\in\nn$ such that $\imageof{\datanum\cdot\empdistribution}\subset\nn$.
%    Then the tuple $(\kg,\impformula,\{\extformulas\})$ defined by a Knowledge Graph
%    \begin{align}
%        \kg =
%        & \bigcup_{\atomindicesin}  \{(
%        s_{j, \atomindices} \quad \mathrm{rdf:type} \quad C ) : j \in [\datanum\cdot\empdistribution(\atomindices)] \}  \\
%        &\bigcup_{\atomindicesin}  \{(
%        s_{j, \atomindices} \quad \mathrm{rdf:type} \quad C_\atomenumerator
%        ) : j \in [\datanum\cdot\empdistribution(\atomindices)], \atomenumeratorin , \atomlegindexof{\atomenumerator}=1\}
%    \end{align}
%    further an importance formula by the query
%    \begin{centeredscript}
%        \impformula = SELECT \{ ?x \} WHERE \{ ?x \quad \rdftype\quad C \, .\}
%    \end{centeredscript}
%    and extraction formulas for each $\atomenumeratorin$ by the query
%    \begin{centeredscript}
%        $\extformulaof{\atomenumerator}$ = SELECT \{ ?x \} WHERE \{ ?x \quad \rdftype \quad $C_\atomenumerator$ \, .\}
%    \end{centeredscript}
%    reproduces $\empdistribution$.
%\end{theorem}
%\begin{proof}
%    With respect to any enumeration of the resources of $\kg$ we have
%    \begin{align}
%        \kggroundingof{\impformula}
%        = \sum_{\atomindicesin} \sum_{j \in [\datanum\cdot\empdistribution(\atomindices)]} \onehotmapof{s_{j, \atomindices} }
%    \end{align}
%    and
%    \begin{align}
%        \kggroundingof{\extformulaof{\atomenumerator}}
%        = \sum_{\atomindicesin \, : \, \atomlegindexof{\atomenumerator} = 1} \sum_{j \in [\datanum\cdot\empdistribution(\atomindices)]} \onehotmapof{s_{j, \atomindices} } \, .
%    \end{align}
%    Summing over the resource variables of these tensors in a contraction we get
%    \begin{align}
%        \contractionof{\{\kggroundingof{\impformula}\}\cup\{\bencodingof{\kggroundingof{\extformulaof{\atomenumerator}}\, : \, \atomenumeratorin}\}}{\shortcatvariables}
%        & = \sum_{\atomenumeratorin}  \datanum\cdot\empdistribution(\atomindices) \cdot \onehotmapof{\atomindices} = \datanum \cdot \empdistribution
%    \end{align}
%    and therefore
%    \begin{align}
%        \normalizationof{\{\kggroundingof{\impformula}\}\cup\{\bencodingof{\kggroundingof{\extformulaof{\atomenumerator}}}\, : \, \atomenumeratorin\}}{\shortcatvariables} = \empdistribution \, .
%    \end{align}
%\end{proof}







\subsect{Samples by Tuples of Objects}

%\paragraph{TBox:} The categorical variables of the factored system are the classes.
%We define atomic formulas by the state indicators of each categorical variable as in \secref{sec:categoricalTN}.
%Each such atomic formula corresponds with a sub-class of the classes.
%By definition, each collection of state indicators define thus pairwise disjoint subclasses.
%
%\paragraph{ABox:} The samples are represented by single individuals in the Knowledge Graph.
%Their sub-class memberships corresponding with the categorical variables of the system are instantiated whenever the atom is true in the sample.
%%\subsubsect{Samples by pairs of resources}
%
%\begin{remark}[Refinement of the Samples]
%    We can split each sample node into a pair of individuals.
%    For this we need to specify, which each class membership will be encoded in a unary or binary attribute of the splitted individuals.
%    This specification is possible based on the extraction query and the atomic formulas.
%\end{remark}
%
%%
%Taking any importance query $\impformula$, which has no permutation symmetries, we can instantiate each projection variable for each sample and prepare the links according to the triple patterns.
%When the atom queries $\extformulas$ have different triple patterns compared with $\impformula$, we instantiate those in cases where $\atomlegindexof{\atomenumerator}=1$.


%
We now instantiate multiple objects for each datapoint, one for each variable of the importance formula, i.e. $\worlddomain=[\datdim]\times[\indorder]$
Label individuals $s_{\datindex,\indenumerator}$ by data index and variable index.

\begin{lemma}
    Let there a data map $\datamap$, queries $\impformula,\shortextformulas$ and a \firstOrderLogic{} world containing objects $s_{\datindex,\indenumerator}$ for $\datindexin$ and $\indenumeratorin$
    If
    \begin{align*}
        \kggroundingof{\impformula}
        = \sum_{\datindexin} \bigotimes_{\indenumeratorin} \onehotmapofat{\indexinterpretationof{s_{\datindex,\indenumerator}}}{\indvariableof{\indenumerator}}
    \end{align*}
    and for any $\atomenumeratorin$
    \begin{align*}
        \kggroundingof{\extformulaof{\atomenumerator}}
        = \sum_{\datindex : \datamap_{\atomenumerator}(\datindex)=1} \bigotimes_{\indvariableof{\indenumerator} \in \indvariableof{\extformulaof{\atomenumerator}}}
        \onehotmapofat{\indexinterpretationof{s_{\datindex,\indenumerator}}}{\indvariableof{\indenumerator}} \, .
    \end{align*}
%    \[ \kggroundingof{\extformulaof{\atomenumerator}}
%    = \sum_{\datindex : \datamap^{\atomenumerator}(\datindex)=1} \bigotimes_{\indenumerator \in \extformulaof{\atomenumerator}} \onehotmapof{\datindex,\indenumerator} \, . \]
    Then the tuple $(\kg,\impformula,\{\extformulas\})$ reproduces $\empdistribution$.
\end{lemma}
\begin{proof}
    We notice that the grounding of the importance formula is in a basis CP format, since by assumption
    \begin{align*}
        \kggroundingof{\impformula}
        = \sum_{\datindexin} \bigotimes_{\indenumeratorin} \onehotmapofat{\indexinterpretationof{s_{\datindex,\indenumerator}}}{\indvariableof{\indenumerator}} \, .
    \end{align*}
    We choose $\datvariable$ to enumerate the non-vanishing entries and get a term selecting map
    \begin{align*}
        \secdatamap_{\indenumerator}(\datindex) = \indexinterpretationof{s_{\datindex,\indenumerator}} \, .
    \end{align*}
    From this we have
    \begin{align*}
        \contractionof{
            \{\bencodingofat{\kggroundingof{\extformulaof{\atomenumerator}}}{\catvariableof{\atomenumerator},\indvariableof{\extformulaof{\atomenumerator}}}\} \cup
            \{\bencodingofat{\secdatamap_{\indenumerator}}{\indvariableof{\indenumerator},\datvariable} \, : \, \indenumeratorin\}
        }{\catvariableof{\catenumerator},\datvariable}
        = \bencodingofat{\datamap_{\atomenumerator}}{\catvariableof{\catenumerator},\datvariable}
    \end{align*}
    and the claim follows with the same argumentation as in the proof of \theref{the:reproducingSingleObjects}.
\end{proof}


%Let us construct a Knowledge Graph
%\begin{align*}
%        \kggroundingof{\rdf}
%        =
%        \sum_{\datindexin}\sum_{\indenumeratorin}
%        \onehotmapof{\indexinterpretationof{s_{\datindex,\indenumerator}}}{\sindvariable}
%        \otimes \onehotmapof{\indexinterpretationof{\mathrdftype}}{\pindvariable}
%        \otimes \onehotmapof{\indexinterpretationof{C}}{\oindvariable}
%        +
%        \sum_{\datindexin} \sum_{\atomenumeratorin \, : \, \datamap_{\atomenumerator}(\datindex)=1} \sum_{\indvariableof{\indenumerator}\in\indvariableof{}}
%        \onehotmapof{\indexinterpretationof{s_{\datindex,\indenumerator}}}{\sindvariable}
%        \otimes \onehotmapof{\indexinterpretationof{\mathrdftype}}{\pindvariable}
%        \otimes \onehotmapof{\indexinterpretationof{C_\atomenumerator}}{\oindvariable} \, .
%\end{align*}
%    We further define an importance formula by the $\sparql$ query
%\begin{centeredscript}
%        \impformula = SELECT \{ ?x_0 \cdots ?x_{\indorder-1} \} WHERE \{ ?x_0 \quad \rdftype\quad C \, .\}
%\end{centeredscript}
%    and for each $\atomenumeratorin$ an extraction formula by the query
%\begin{centeredscript}
%        $\extformulaof{\atomenumerator}$ = SELECT \{ ?x \} WHERE \{ ?x \quad \rdftype \quad $C_\atomenumerator$ \, .\} \, .
%\end{centeredscript}
%    Then the triple $(\kg,\impformula,\shortextformulas)$ reproduces $\empdistribution$.



\sect{Discussion}


% Probabilistic Relational Models
Statistical Models are called Probabilistic Relational Models. % (RUSSELL - Chapter Probabilistic Programming).
Extensions are models that also handle structural uncertainty, i.e. distributions of worlds with varying $\worlddomain$.

% Comparison with network science
In the emerging area of network science \cite{barabasi_network_2016, giovanni_russo_vito_latora_complex_2017}, statistical models for random graphs are investigated.
Statistical Models of \firstOrderLogic{} go beyond the typical single edge type perspective of network science.


%
\begin{remark}[Alternative Representation of empirical distributions]
    So far, we have motivated the representation of empirical distributions based on basis CP decompositions based on data maps.
    In this section, based on the extraction queries, we have observed that empirical distributions might have more efficient representation formats.
    In many applications such as the computation of log-likelihoods we can use any representation of the empirical distribution by tensor networks.
    It is thus not necessary to compute the data cores as above, unless one requires a list of the extracted samples.
\end{remark}